% Title: Lagrangian Smoothability of Quadrivalent Polyhedral Surfaces

%%% PROBLEM STATEMENT %%%

\noindent\textbf{Problem.}
A polyhedral Lagrangian surface $K$ in $\mathbb{R}^4$ is a finite polyhedral
complex all of whose faces are Lagrangians, and which is a topological
submanifold of $\mathbb{R}^4$. A Lagrangian smoothing of $K$ is a
Hamiltonian isotopy $K_t$ of smooth Lagrangian submanifolds, parameterised
by $(0,1]$, extending to a topological isotopy, parametrised by $[0,1]$,
with endpoint $K_0 = K$. Let $K$ be a polyhedral Lagrangian surface with
the property that exactly $4$ faces meet at every vertex. Does $K$
necessarily have a Lagrangian smoothing?

%%% SOLUTION %%%

\bigskip
\noindent\textbf{Answer.} \textit{Yes.} Every quadrivalent polyhedral
Lagrangian surface admits a Lagrangian smoothing. We construct the smoothing
via the Lagrangian $h$-principle (Gromov--Lees--Eliashberg), using a Weinstein
tubular neighborhood to reduce the problem to smoothing an exact Lagrangian
immersion that is already an embedding.

\medskip
\noindent\textbf{Strategy.}
The proof has five stages:
(1)~show that $K$ admits a Weinstein tubular neighborhood modeled on
$T^*K_{\mathrm{smooth}}$;
(2)~construct a smooth immersion $\iota_t: \Sigma\to\mathbb{R}^4$ that is
$C^0$-close to $K$ and is a Lagrangian immersion for each $t\in(0,1]$;
(3)~apply the Gromov--Lees $h$-principle to promote the immersion to an
exact Lagrangian;
(4)~use the 4-valent (generic) condition to ensure the immersion is
actually an embedding;
(5)~produce the Hamiltonian isotopy via the Moser trick.

% =========================================================================
\section*{1.\ Setup and notation}
\label{sec:p8-setup}
% =========================================================================

We work in $(\mathbb{R}^4, \omega_0)$ with coordinates $(x_1,y_1,x_2,y_2)$
and symplectic form $\omega_0 = dx_1\wedge dy_1 + dx_2\wedge dy_2$. The
Liouville one-form is $\lambda_0 = y_1\,dx_1 + y_2\,dx_2$, so
$\omega_0 = -d\lambda_0$.

Let $K\subset\mathbb{R}^4$ be a quadrivalent polyhedral Lagrangian surface.
Write $\Sigma$ for the underlying abstract topological surface (compact,
without boundary), and $\phi: \Sigma\to\mathbb{R}^4$ for the piecewise-linear
embedding with image $K$.

\begin{definition}
The \emph{singular set} of $K$ is $K_{\mathrm{sing}} = \mathcal{E}\cup\mathcal{V}$,
where $\mathcal{E}$ is the (closed) edge locus and
$\mathcal{V}=\{v_1,\ldots,v_N\}$ is the vertex set. The \emph{smooth part}
$K_{\mathrm{reg}} = K\setminus K_{\mathrm{sing}}$ is a smooth Lagrangian
submanifold (each face is an open convex polygon in a Lagrangian plane).
\end{definition}

\begin{definition}[Lagrangian Grassmannian]
$\Lambda(2)$ denotes the Grassmannian of Lagrangian $2$-planes in
$(\mathbb{R}^4,\omega_0)$. It is diffeomorphic to $U(2)/O(2)$ and satisfies
$\pi_1(\Lambda(2))\cong\mathbb{Z}$ (generated by the Maslov class).
\end{definition}

% =========================================================================
\section*{2.\ The Gauss map and its extension}
\label{sec:p8-gauss}
% =========================================================================

On the smooth part $K_{\mathrm{reg}}$, the tangent plane at each point is
a Lagrangian plane, defining the \emph{Gauss map}
\[
  G: K_{\mathrm{reg}} \;\to\; \Lambda(2), \qquad x\;\mapsto\; T_xK.
\]
This map is locally constant on each face.

\begin{proposition}[Maslov index vanishing at vertices]
\label{prop:p8-maslov}
Let $v\in\mathcal{V}$ be a quadrivalent vertex with cyclically ordered
Lagrangian tangent planes $\Pi_1,\Pi_2,\Pi_3,\Pi_4$. Consecutive planes
share an edge direction: $\Pi_i\cap\Pi_{i+1}$ contains a common line
$\ell_i$ (indices mod~$4$). The Maslov index of the loop
$\gamma_v: \Pi_1\to\Pi_2\to\Pi_3\to\Pi_4\to\Pi_1$ in $\Lambda(2)$
vanishes: $\mu(\gamma_v)=0$.
\end{proposition}

\begin{proof}
Identify $\mathbb{R}^4\cong\mathbb{C}^2$ via $z_j = x_j + iy_j$.
Each Lagrangian plane through the origin is $U\cdot\mathbb{R}^2$ for some
$U\in U(2)$, defined up to right multiplication by $O(2)$. The Maslov index
of a loop is the winding number of
$(\det U)^2: S^1\to S^1$.

At each transition $\Pi_i\to\Pi_{i+1}$, the two planes share the line
$\ell_i$. In a unitary frame adapted to $\ell_i$, the transition replaces
one real basis vector by another while fixing $\ell_i$. This corresponds to
an element $T_i\in O(2)$ acting on the right, with $\det T_i = -1$
(a reflection). Thus $\det U_{i+1} = \det U_i\cdot\det T_i = -\det U_i$.

After four transitions:
$\det U_1\to -\det U_1\to \det U_1\to -\det U_1\to \det U_1$.
The squared determinant returns to its original value without winding:
$(\det U)^2$ is constant at each vertex of the piecewise loop. Along
each geodesic arc in $\Lambda(2)$ connecting consecutive planes, the
squared determinant traces a path that starts and ends at the same value
(since both endpoints give the same $(\det U)^2$ after two steps). Summing
over all four arcs, the total winding is zero.
\end{proof}

\begin{lemma}[Extension of the Gauss map]
\label{lem:p8-gauss-ext}
The Gauss map $G: K_{\mathrm{reg}}\to\Lambda(2)$ extends to a continuous map
$\overline{G}: \Sigma\to\Lambda(2)$ (where $\Sigma$ is the abstract surface).
\end{lemma}

\begin{proof}
At each edge point $p\in\mathcal{E}$, the two adjacent Lagrangian planes
share a line, so the geodesic midpoint in $\Lambda(2)$ provides a canonical
extension. At each vertex, Proposition~\ref{prop:p8-maslov} shows
$\mu(\gamma_v)=0$, so the loop bounds a disk in $\Lambda(2)$
($\pi_1(\Lambda(2))\cong\mathbb{Z}$ and the Maslov index is the winding
number), which provides the extension.
\end{proof}

% =========================================================================
\section*{3.\ Formal Lagrangian data: the $h$-principle setup}
\label{sec:p8-hprinciple}
% =========================================================================

\begin{definition}[Formal Lagrangian immersion]
A \emph{formal Lagrangian immersion} of a surface $\Sigma$ into
$(\mathbb{R}^4,\omega_0)$ is a pair $(f, F_s)$ where:
\begin{itemize}
\item $f: \Sigma\to\mathbb{R}^4$ is a smooth map (not necessarily an immersion);
\item $F_s: T\Sigma\to f^*T\mathbb{R}^4$, $s\in[0,1]$, is a homotopy of
bundle monomorphisms with $F_0 = df$ (when $f$ is an immersion) and
$F_1(\Sigma)\subset\Lambda(2)$ (i.e., $F_1$ maps each tangent plane to a
Lagrangian plane).
\end{itemize}
\end{definition}

\begin{theorem}[Gromov--Lees {\cite{Gromov86, Lees76}}]
\label{thm:p8-gromov-lees}
Every formal Lagrangian immersion of a closed surface~$\Sigma$ into
$(\mathbb{R}^4,\omega_0)$ is homotopic (through formal Lagrangian immersions)
to a genuine Lagrangian immersion. Moreover, the genuine Lagrangian immersion
can be chosen $C^0$-close to the original map~$f$.
\end{theorem}

Our polyhedral Lagrangian $K = \phi(\Sigma)$ naturally provides formal
Lagrangian data:
\begin{itemize}
\item The map $f := \phi: \Sigma\to\mathbb{R}^4$ is a PL embedding.
\item The extended Gauss map $\overline{G}: \Sigma\to\Lambda(2)$ from
Lemma~\ref{lem:p8-gauss-ext} provides the Lagrangian tangent plane data.
\end{itemize}

After a small perturbation of $\phi$ to a smooth immersion $f_\epsilon$
(staying $C^0$-close to $\phi$), the pair $(f_\epsilon, \overline{G})$ is
a formal Lagrangian immersion. Theorem~\ref{thm:p8-gromov-lees} promotes this
to a genuine smooth Lagrangian immersion $\iota: \Sigma\looparrowright
\mathbb{R}^4$ that is $C^0$-close to~$K$.

% =========================================================================
\section*{4.\ From immersion to embedding: the 4-valent condition}
\label{sec:p8-embedding}
% =========================================================================

The $h$-principle gives an immersion, not an embedding. We now show that
for quadrivalent polyhedral Lagrangians, the immersion can be taken to be
an embedding.

\begin{proposition}[Genericity of the 4-valent condition]
\label{prop:p8-generic}
The condition that exactly $4$ faces meet at every vertex is the
\emph{generic} condition for polyhedral Lagrangian surfaces in
$\mathbb{R}^4$. Equivalently, at each vertex, the four edge directions span
a $4$-dimensional subspace of $\mathbb{R}^4$ (general position).
\end{proposition}

\begin{proof}
At a vertex $v$, the four edge directions $\ell_1,\ell_2,\ell_3,\ell_4$
are lines in $\mathbb{R}^4$. By the normal form of
Section~\ref{sec:p8-vertex-normal} below, after a symplectomorphism we can
arrange $\ell_1\parallel\partial_{x_1}$,
$\ell_2\parallel\partial_{x_2}$,
$\ell_3\parallel\partial_{x_1}+\beta\partial_{y_1}$ ($\beta\ne 0$), and
$\ell_4\parallel\partial_{x_2}+\delta\partial_{y_2}$ ($\delta\ne 0$). These
four directions span all of $\mathbb{R}^4$.
\end{proof}

\begin{lemma}[Vertex normal form]
\label{sec:p8-vertex-normal}
At a quadrivalent vertex $v$ (placed at the origin), a linear
symplectomorphism brings the four Lagrangian planes to:
\begin{align*}
\Pi_1 &= \mathrm{span}(\partial_{x_2}+\delta\partial_{y_2},\;
          \partial_{x_1}), \\
\Pi_2 &= \mathrm{span}(\partial_{x_1},\;\partial_{x_2}), \\
\Pi_3 &= \mathrm{span}(\partial_{x_2},\;
          \partial_{x_1}+\beta\partial_{y_1}), \\
\Pi_4 &= \mathrm{span}(\partial_{x_1}+\beta\partial_{y_1},\;
          \partial_{x_2}+\delta\partial_{y_2}),
\end{align*}
with $\beta,\delta\ne 0$. All four planes are Lagrangian (verified by
$\omega_0(u,w)=0$ for each pair of spanning vectors).
\end{lemma}

\begin{theorem}[Embedding from the $h$-principle]
\label{thm:p8-embedding}
Let $\iota: \Sigma\looparrowright\mathbb{R}^4$ be a Lagrangian immersion
$C^0$-close to the PL embedding $\phi: \Sigma\to\mathbb{R}^4$. If the
$C^0$-closeness is sufficient (controlled by the injectivity radius of the
embedding $\phi$), then $\iota$ is an embedding.
\end{theorem}

\begin{proof}
Since $\phi$ is a topological embedding of a compact surface $\Sigma$, it
has a positive \emph{normal injectivity radius}
$\rho := \inf_{p\ne q\in\Sigma} |\phi(p)-\phi(q)| / d_\Sigma(p,q) > 0$
(where $d_\Sigma$ is the intrinsic distance on $\Sigma$ with some fixed
Riemannian metric). For $||\iota - \phi||_{C^0} < \rho/3$, the map $\iota$
is injective by the triangle inequality: if $\iota(p) = \iota(q)$ with
$p\ne q$, then
$|\phi(p)-\phi(q)| \le |\phi(p)-\iota(p)| + |\iota(p)-\iota(q)| +
|\iota(q)-\phi(q)| < 2\rho/3$,
but $|\phi(p)-\phi(q)| \ge \rho\cdot d_\Sigma(p,q) > 0$, a contradiction
when $d_\Sigma(p,q)$ is bounded below.

More carefully: $\phi$ is a homeomorphism onto its image, so for any
$\epsilon>0$, there exists $\delta>0$ such that
$d_\Sigma(p,q)\ge\epsilon$ implies $|\phi(p)-\phi(q)|\ge\delta$. For
$p,q$ with $d_\Sigma(p,q)<\epsilon$, the immersion $\iota$ is an embedding
on $\epsilon$-balls (since it is a smooth immersion close to a PL embedding
that is locally injective). Combining these two cases yields global
injectivity.

The 4-valent condition (Proposition~\ref{prop:p8-generic}) ensures that the
4 edge directions span $\mathbb{R}^4$, which means the PL embedding $\phi$
has no ``self-tangencies'' at vertices (distinct sheets of $K$ separate in
all four coordinate directions). This provides the positive normal
injectivity radius.
\end{proof}

% =========================================================================
\section*{5.\ Exactness and the Hamiltonian isotopy (Moser trick)}
\label{sec:p8-moser}
% =========================================================================

At this stage we have a smooth Lagrangian embedding
$\iota: \Sigma\hookrightarrow\mathbb{R}^4$ that is $C^0$-close to~$K$.
We now construct the full Hamiltonian isotopy.

\begin{proposition}[Exactness]
\label{prop:p8-exact}
The Lagrangian embedding $\iota(\Sigma)$ is \emph{exact}: the Liouville
form $\lambda_0$ restricts to an exact $1$-form on $\iota(\Sigma)$. That is,
$\iota^*\lambda_0 = dh$ for some smooth function $h: \Sigma\to\mathbb{R}$.
\end{proposition}

\begin{proof}
In $(\mathbb{R}^4, \omega_0 = d\lambda_0)$, every compact Lagrangian
submanifold is exact. Indeed, the symplectic form on $\mathbb{R}^4$ is
exact ($\omega_0 = -d\lambda_0$), so for any Lagrangian $L$:
$\iota^*\omega_0 = 0$ implies $d(\iota^*\lambda_0) = 0$, i.e.,
$\iota^*\lambda_0$ is closed. Since $\Sigma$ is compact and
$H^1_{\mathrm{dR}}(\Sigma;\mathbb{R})$ is finite-dimensional, we need to
check that $[\iota^*\lambda_0]=0\in H^1(\Sigma;\mathbb{R})$.

For the Lagrangian $\iota(\Sigma)$ that is $C^0$-close to the PL surface~$K$:
each face of $K$ lies in a Lagrangian \emph{plane} through some point, and
on a Lagrangian plane $\Pi$ with base point $p$, the restriction
$\lambda_0|_\Pi = d(\text{quadratic function})$ is exact. The smooth
Lagrangian $\iota(\Sigma)$ is obtained by the $h$-principle as a small
perturbation, and the cohomology class $[\iota^*\lambda_0]$ varies
continuously. Since it starts at zero (the PL surface has exact restriction
on each face, and the global class is zero by telescoping around any cycle),
it remains zero.
\end{proof}

\begin{theorem}[Hamiltonian isotopy via Moser]
\label{thm:p8-moser}
There exists a smooth family of Lagrangian embeddings
$K_t: \Sigma\hookrightarrow\mathbb{R}^4$, $t\in(0,1]$, such that:
\begin{enumerate}
\item[(a)] Each $K_t(\Sigma)$ is a smooth exact Lagrangian submanifold.
\item[(b)] As $t\to 0^+$, $K_t(\Sigma)\to K$ in the Hausdorff metric.
\item[(c)] The family extends to a continuous family for $t\in[0,1]$ with
$K_0 = K$.
\item[(d)] The isotopy $\{K_t\}_{t\in(0,1]}$ is Hamiltonian: there exists
a smooth family of Hamiltonian functions $H_t: \mathbb{R}^4\to\mathbb{R}$
such that $\frac{d}{dt}K_t = X_{H_t}\circ K_t$, where $X_{H_t}$ is the
Hamiltonian vector field of $H_t$.
\end{enumerate}
\end{theorem}

\begin{proof}
\textbf{Step 1: Parametrized smoothing.}
For each $t\in(0,1]$, define a smoothing parameter $\epsilon(t) = t$ and
apply the Gromov--Lees construction with smoothing at scale $\epsilon(t)$.
Concretely, smooth the corners of $K$ within an $\epsilon(t)$-neighborhood
of the singular set $K_{\mathrm{sing}}$. By Sections~3--4, the result is
a smooth exact Lagrangian embedding $K_t$ for each $t>0$, and
$K_t\to K$ in $C^0$ as $t\to 0$. This gives~(a)--(c).

\textbf{Step 2: Hamiltonian isotopy.}
For $0 < s < t \le 1$, the Lagrangians $K_s(\Sigma)$ and $K_t(\Sigma)$
are both exact Lagrangian embeddings that are Lagrangian isotopic (through
$\{K_r\}_{r\in[s,t]}$). We promote this to a \emph{Hamiltonian} isotopy
using the Moser method.

Consider the family $\omega_r := K_r^*\omega_0 = 0$ (each $K_r$ is
Lagrangian, so the pullback vanishes). The key is to construct the
Hamiltonian function generating the isotopy.

By Proposition~\ref{prop:p8-exact}, for each $r$, there exists
$h_r: \Sigma\to\mathbb{R}$ with $K_r^*\lambda_0 = dh_r$. The family $h_r$
can be chosen to depend smoothly on~$r$ (by the smooth dependence of the
$h$-principle construction on parameters).

Define the Hamiltonian $H_r: \mathbb{R}^4\to\mathbb{R}$ by extending $h_r$
from the Lagrangian $K_r(\Sigma)$ to a neighborhood using a Weinstein tubular
neighborhood. Explicitly: by Weinstein's Lagrangian neighborhood theorem
\cite{Weinstein71}, a neighborhood of $K_r(\Sigma)$ in $\mathbb{R}^4$ is
symplectomorphic to a neighborhood of the zero section in
$T^*K_r(\Sigma)$. In cotangent coordinates $(q,p)$, the Lagrangian
$K_r(\Sigma)$ is the zero section $\{p=0\}$, and the primitive
$\lambda_0|_{K_r} = dh_r$ corresponds to $h_r$ on the zero section.
Extend $H_r$ by pullback: $H_r(q,p) = h_r(q)$ (constant in the fiber
direction).

The Hamiltonian vector field $X_{H_r}$ on the cotangent bundle satisfies
$\iota_{X_{H_r}}\omega = -dH_r$. On the zero section, $X_{H_r}$ is tangent
to the zero section and generates the desired flow. By construction, this
flow maps $K_r(\Sigma)$ to $K_{r+\delta}(\Sigma)$ for small $\delta$,
giving~(d).

\textbf{Step 3: Global Hamiltonian.}
The Hamiltonian $H_t$ can be extended from the tubular neighborhood to all
of $\mathbb{R}^4$ using a smooth cutoff function (multiplying by a bump
function supported in the tubular neighborhood). This does not affect the
flow on the Lagrangian, since $X_{H_t}$ restricted to $K_t(\Sigma)$
depends only on $H_t|_{K_t(\Sigma)}$. Since $\mathbb{R}^4$ is contractible,
there is no topological obstruction to this extension.
\end{proof}

% =========================================================================
\section*{6.\ Detailed verification of the local models}
\label{sec:p8-local}
% =========================================================================

\subsection*{6.1.\ Edge smoothing via generating functions}

\begin{lemma}[Edge normal form and smoothing]
\label{lem:p8-edge-smooth}
At an edge where two faces $F_1, F_2$ meet, choose Darboux coordinates
with the edge along the $x_1$-axis:
\begin{align*}
F_1 &= \{y_1=0,\, y_2=0,\, x_2\ge 0\}, \\
F_2 &= \{y_1=0,\, y_2=\alpha x_2,\, x_2\le 0\} \quad (\alpha\ne 0).
\end{align*}
The generating function $S(x_1,x_2) = \alpha\cdot h_\epsilon(x_2)$, where
$h_\epsilon$ is a smooth convex function with $h_\epsilon''(x_2)=1$ for
$x_2\le -\epsilon$ and $h_\epsilon''(x_2)=0$ for $x_2\ge\epsilon$, produces
the smooth Lagrangian surface
\[
  L_\epsilon = \{(x_1,\, 0,\, x_2,\, \alpha h_\epsilon'(x_2))\}
\]
that interpolates between the two faces outside the $\epsilon$-strip
$|x_2|<\epsilon$.
\end{lemma}

\begin{proof}
The surface $L_\epsilon = \mathrm{graph}(dS) = \{(x, \partial_x S(x))\}$
is automatically Lagrangian in $T^*\mathbb{R}^2 = \mathbb{R}^4$, since
$\omega_0 = \sum dx_i\wedge dy_i$ and $y_i = \partial_{x_i} S$ implies
\[
  \omega_0|_{L_\epsilon}
  = \sum_i d(\partial_{x_i}S)\wedge dx_i
  = \sum_{i,j} \frac{\partial^2 S}{\partial x_j\partial x_i}\,dx_j\wedge dx_i
  = 0
\]
by the symmetry of the Hessian ($\partial^2 S/\partial x_j\partial x_i =
\partial^2 S/\partial x_i\partial x_j$).

For $x_2\ge\epsilon$: $h_\epsilon'(x_2)=0$, so $L_\epsilon = \{y_1=0,y_2=0\}
= F_1$. For $x_2\le -\epsilon$: $h_\epsilon'(x_2) = x_2$ (since
$h_\epsilon'' = 1$ there), so $L_\epsilon = \{y_1=0, y_2=\alpha x_2\} = F_2$.
\end{proof}

\subsection*{6.2.\ Vertex smoothing via generating functions}

\begin{lemma}[Vertex smoothing]
\label{lem:p8-vertex-smooth}
At a quadrivalent vertex in the normal form of Lemma~\ref{sec:p8-vertex-normal},
the generating function
\[
  S(x_1,x_2)
  \;=\;
  \beta\cdot\frac{x_1^2}{2}\cdot\chi_\epsilon(x_2)
  \;+\;
  \delta\cdot\frac{x_2^2}{2}\cdot\chi_\epsilon(x_1),
\]
where $\chi_\epsilon:\mathbb{R}\to[0,1]$ is a smooth non-increasing cutoff
with $\chi_\epsilon(u)=1$ for $u\le -\epsilon$ and $\chi_\epsilon(u)=0$ for
$u\ge\epsilon$, produces a smooth Lagrangian surface that matches all four
faces outside a neighborhood of the vertex.
\end{lemma}

\begin{proof}
The surface $L_\epsilon^v = \mathrm{graph}(dS)$ is Lagrangian by the same
Hessian-symmetry argument as Lemma~\ref{lem:p8-edge-smooth}: the graph of
the gradient of any smooth function $S(x_1,x_2)$ is automatically Lagrangian.

We compute the gradients:
\begin{align*}
  y_1 = \partial_{x_1}S
  &= \beta\, x_1\,\chi_\epsilon(x_2)
     + \delta\,\tfrac{x_2^2}{2}\,\chi_\epsilon'(x_1), \\
  y_2 = \partial_{x_2}S
  &= \beta\,\tfrac{x_1^2}{2}\,\chi_\epsilon'(x_2)
     + \delta\, x_2\,\chi_\epsilon(x_1).
\end{align*}

We verify boundary matching in each of the four quadrants (where the
cutoff derivatives vanish):

\textbf{Quadrant I} ($x_1\gg\epsilon$, $x_2\gg\epsilon$):
$\chi_\epsilon(x_1)=\chi_\epsilon(x_2)=0$, derivatives also vanish, so
$y_1=y_2=0$. This matches $\Pi_2=\{y_1=0,\,y_2=0\}$.

\textbf{Quadrant II} ($x_1\ll -\epsilon$, $x_2\gg\epsilon$):
$\chi_\epsilon(x_1)=1$, $\chi_\epsilon(x_2)=0$, derivatives vanish. Then
$y_1 = 0$, $y_2 = \delta\, x_2$. This matches
$\Pi_1 = \{y_1=0,\, y_2=\delta x_2\}$.

\textbf{Quadrant III} ($x_1\gg\epsilon$, $x_2\ll -\epsilon$):
$\chi_\epsilon(x_1)=0$, $\chi_\epsilon(x_2)=1$, derivatives vanish. Then
$y_1 = \beta\, x_1$, $y_2 = 0$. This matches
$\Pi_3=\{y_1=\beta x_1,\, y_2=0\}$.

\textbf{Quadrant IV} ($x_1\ll -\epsilon$, $x_2\ll -\epsilon$):
$\chi_\epsilon(x_1)=\chi_\epsilon(x_2)=1$, derivatives vanish. Then
$y_1 = \beta\, x_1$, $y_2 = \delta\, x_2$. This matches
$\Pi_4=\{y_1=\beta x_1,\, y_2=\delta x_2\}$.
\end{proof}

\subsection*{6.3.\ Compatibility of local models}

\begin{proposition}[Global patching]
\label{prop:p8-patch}
The local edge and vertex smoothing models are compatible on their overlaps
and assemble into a global smooth Lagrangian surface $K_t$ for each
$t=\epsilon\in(0,1]$.
\end{proposition}

\begin{proof}
The edge smoothing (Lemma~\ref{lem:p8-edge-smooth}) modifies $K$ only within
an $\epsilon$-strip around each edge. The vertex smoothing
(Lemma~\ref{lem:p8-vertex-smooth}) modifies $K$ only within an
$\epsilon$-ball around each vertex. By choosing $\epsilon$ smaller than half
the minimum edge length and half the minimum vertex-to-vertex distance, the
vertex neighborhoods are disjoint from each other and from the edge-only
smoothing regions.

In the overlap region (near an edge but within a vertex neighborhood), the
vertex generating function $S^v$ already produces the correct edge smoothing
on each edge emanating from that vertex. This is because the generating
function $S^v$ was designed to match the face data along each edge direction
(Lemma~\ref{lem:p8-vertex-smooth}). By a partition-of-unity interpolation
between $S^v$ (near the vertex) and $S^e$ (away from the vertex along the
edge), the two models agree in the overlap.

The Lagrangian condition is preserved because the interpolation is performed
at the level of generating functions (not at the level of surfaces), and
the graph of the gradient of any smooth function is Lagrangian.
\end{proof}

% =========================================================================
\section*{7.\ Main theorem}
\label{sec:p8-main}
% =========================================================================

\begin{theorem}[Main result]
\label{thm:p8-main}
Let $K\subset\mathbb{R}^4$ be a polyhedral Lagrangian surface with exactly
$4$ faces meeting at every vertex. Then $K$ admits a Lagrangian smoothing:
there exists a Hamiltonian isotopy $\{K_t\}_{t\in(0,1]}$ of smooth
Lagrangian submanifolds extending to a topological isotopy on $[0,1]$ with
$K_0 = K$.
\end{theorem}

\begin{proof}
Combine the preceding results:
\begin{enumerate}
\item By Proposition~\ref{prop:p8-maslov} and Lemma~\ref{lem:p8-gauss-ext},
the Gauss map extends over $\Sigma$, providing formal Lagrangian data.
\item By the Gromov--Lees $h$-principle (Theorem~\ref{thm:p8-gromov-lees}),
the formal data is realized by a genuine Lagrangian immersion.
\item By Theorem~\ref{thm:p8-embedding}, the 4-valent condition ensures the
immersion is an embedding for small enough smoothing parameter.
\item By Proposition~\ref{prop:p8-exact}, the embedding is exact Lagrangian.
\item By Theorem~\ref{thm:p8-moser}, the family assembles into a Hamiltonian
isotopy.
\item By Proposition~\ref{prop:p8-patch}, the local smoothing models
(Lemmas~\ref{lem:p8-edge-smooth} and~\ref{lem:p8-vertex-smooth}) are globally
compatible.
\end{enumerate}
\end{proof}

\begin{remark}[Role of the 4-valent condition]
The restriction to exactly $4$ faces per vertex is crucial for two reasons:
(a)~it is the generic valence for a polyhedral surface in $\mathbb{R}^4$
(four edges generically span the ambient $4$-dimensional space), ensuring
the normal injectivity radius is positive; and
(b)~at each vertex, the four Lagrangian planes are arranged so that the Maslov
index vanishes (Proposition~\ref{prop:p8-maslov}), which is necessary for
the Gauss map to extend. For higher valence ($\ge 6$), the Maslov index can
be nonzero, obstructing the extension.
\end{remark}

\begin{remark}[Comparison with other approaches]
This proof avoids the conormal fibration approach of Abouzaid and the
direct local-to-global gluing approach of OpenAI. Instead, it leverages the
$h$-principle for Lagrangian immersions (Gromov--Lees), which automatically
handles the global compatibility that is the main difficulty in direct
constructions. The 4-valent condition enters via the embedding step
(Theorem~\ref{thm:p8-embedding}), not via the Lagrangian condition itself.
\end{remark}

\paragraph{Confidence.} MEDIUM --- The proof strategy is sound: the
$h$-principle for Lagrangian immersions is a well-established deep theorem,
and the reduction from polyhedral to smooth via generating functions is
standard in symplectic topology. The main subtlety is the promotion from
immersion to embedding (Theorem~\ref{thm:p8-embedding}), which relies on
the $C^0$-closeness provided by the $h$-principle and the positive normal
injectivity radius guaranteed by the 4-valent condition. While the argument
is morally correct, a fully rigorous treatment would require careful
estimates on the $C^0$-norm from the $h$-principle (which Gromov's theory
provides in principle but which require nontrivial work to make explicit).
The generating-function local models (Section~6) are verified numerically
in the companion script.

%%% REFERENCES %%%

\bibitem{Gromov86} M.\,Gromov.
\emph{Partial Differential Relations}.
Ergebnisse der Mathematik~9, Springer, 1986.

\bibitem{Lees76} J.\,A.\,Lees.
On the classification of Lagrange immersions.
\emph{Duke Math.\ J.}, 43(2):217--224, 1976.

\bibitem{Weinstein71} A.\,Weinstein.
Symplectic manifolds and their Lagrangian submanifolds.
\emph{Adv.\ Math.}, 6:329--346, 1971.

\bibitem{EliashbergMishachev} Y.\,Eliashberg, N.\,Mishachev.
\emph{Introduction to the $h$-Principle}.
Graduate Studies in Mathematics~48, AMS, 2002.

\bibitem{McDuffSalamon} D.\,McDuff, D.\,Salamon.
\emph{Introduction to Symplectic Topology}.
3rd edition, Oxford University Press, 2017.

\bibitem{ArnoldGivental} V.\,I.\,Arnold, A.\,B.\,Givental.
Symplectic geometry.
In \emph{Dynamical Systems IV}, Encyclopaedia of Math.\ Sci.~4,
pp.~1--138, Springer, 2001.

%%% HSEXPLAIN %%%

Imagine a crumpled piece of paper sitting in four-dimensional space. This
paper is special: it is made of flat triangular or rectangular patches, each
lying in a ``Lagrangian plane'' -- a two-dimensional surface that is
invisible to a certain mathematical measuring device called the symplectic
form. Where patches meet along edges, there are creases. Where edges meet,
there are sharp corners. The question is: can you iron out all the creases
and corners to get a perfectly smooth surface, while keeping the surface
invisible to the symplectic form at every stage of the ironing process?

The answer is yes, provided each corner is ``four-valent'' -- exactly four
patches meet at every corner point. This is the generic situation in four
dimensions, analogous to how three edges typically meet at a corner of a
crumpled surface in three-dimensional space.

The proof works in three main steps. First, we examine each corner and
verify that the arrangement of the four patches has zero ``winding number''
(technically, zero Maslov index). This means the patches cycle through
their orientations in a balanced way, returning to where they started
without any twist. This no-twist condition is what allows the smoothing
to proceed.

Second, we invoke a powerful general principle from the 1970s and 1980s
called the ``h-principle'' (developed by Gromov and others). In plain
terms, it says: if you can imagine a smooth surface that approximately
satisfies the Lagrangian condition, and you can also imagine a smooth
field of Lagrangian tangent planes covering your surface, then you can
actually find a genuinely smooth Lagrangian surface nearby. The flat
patches already give us the Lagrangian tangent planes, and the no-twist
condition at corners lets us smoothly interpolate these planes across
the creases and corners.

Third, we check that the resulting smooth surface does not cross itself.
The four-valent condition is crucial here: four edge directions in
four-dimensional space generically point in independent directions, which
means the patches separate cleanly in all dimensions. This separation
guarantees that when we smooth the corners, nearby parts of the surface
do not accidentally overlap.

Finally, we show the ironing process can be done in a ``Hamiltonian'' way,
meaning there is an energy function driving the smoothing at each stage.
This uses a classical technique called the Moser trick, which converts any
smooth family of special surfaces into one driven by an energy function,
provided the surfaces remain invisible to the symplectic form throughout.
